% !TEX root = handbook-physics.tex
% !TEX TS-program = pdflatex
% !TEX encoding = UTF-8 Unicode






\chapter{Free Harmonic Oscillations}
\label{ch-28}

\section{Harmonic Oscillations}
\label{sec-28.1} 


\subsection{}\label{28.1.1} Oscillations, including vibrations and variations, are processes (motions or changes of state) that are repeated more or less regularly with time. According to the physical nature of the oscillatory process and the “mechanism”by which it is initiated, distinction is made between mechanical oscillations, commonlycalledvibrations(swingingofpendulums; vibration of strings; reciprocation and oscillation of machine components; vibration of buildings, bridges and other structures; compres­ sion and rarefaction of air as sound is propagated through it; rolling, pitching and heaving of ships; sea waves; etc.), electro­ magnetic oscillations (oscillations of alternating electric current in a circuit, variations in the field strength and magnetic induc­ tion vectors, E and B, of an alternating electromagnetic field, etc.), and electromechanical vibrations (vibration of the dia­ phragm of a telephone receiver, of the diaphragm in an elec­ tromagnetic loudspeaker, etc.) and others.
The system executing the oscillation (or vibration) is called the oscillatory (vibratory) system. Free (natural) oscillations are ones that occur in the absence of externally applied variable action on the system. They are due to some initial deviation of the system from a state of stable equilibrium. Forced oscil­ lations (vibrations) are ones set up in a system as a result of variable external effects (for instance, current oscillation in an electric circuit caused by a varying emf, swinging of a pendulum due to varying external forces, etc.).

\subsection{}\label{28.1.2} Oscillations (vibrations) are said to be periodic if all the physical quantities that characterize the oscillatory (vibra­ tory) system and vary during its oscillation are repeated in successive time intervals. The minimum time interval T com-

plying with this condition is called the period of oscillation (or vibration). This is the time required for one complete oscil­ lation (vibration) of the system. The frequency of periodic os­ cillations (vibrations) is the quantity v = 1IT, equal to the number of complete oscillations (vibrations) executed by the system in unit time. The angular, or circular, frequency is the quantity to = 2jiv = 2n/T, which is equal to the number of complete oscillations accomplished in 2n units of time.

\subsection{}\label{28.1.3} In periodic oscillations the dependence of a varying quantitysonthetimetcomplieswiththeconditions(t+ T)=
= 5(*)•
The periodic oscillation of the quantity s (t) is said to be har­
monic if
s (t) = A sin (  t + (p0) or s (t) = A cos (tot + Ti)»
where to = 2nv = (2n/T) = const is the angular, or circular, frequency of harmonic vibration, A = smax = const > 0 is the maximum value of the varying quantity s and is called the amplitude of oscillation, and <p0 and Ti = To — jt/2 are con­ stant quantities. The value of s at some arbitrary instant of time t is determined by the phase of the oscillation   (t) = = tot -]- cp0 [and, respectively, (t) =  J + <pj. The quanti­ ties (p0 and cpx are the initial phases of the oscillation, i.e. the values of   (t) and   x (t) at the instant (t = 0) of time refer­ ence: cp0=  -(0) and Ti =   x(0).

\subsection{}\label{28.1.4} The first and second time derivatives of a harmonically varying quantity s (t) also vary harmonically at the same angu­ lar frequency:
ds /
= A cos ( *+ (po)= A  sin ( *H“ToH—2~
  •— — A  2sin ( £+ cp0)= Am2sin ( £-[-cp0+ n)»
where the amplitudes of ds/dt and d2sldt2are equal, respectively, to  4co and Am2. The initial phase of ds/dt is equal to (cp0 + Jt/2), i.e. the phase difference of the variations of ds/dt ana s is con­ stant and equal to jx/2 (the quantity ds/dt has a phase lead of jt/2 with respect to s). The initial phase of d2s/dt2As equal, tq (To + n)» he. the phase difference of the variations of d2s/dt2

and s is constant and equal to n (the quantity d2sldt2 has a phase lead of jx with respect to s). Curves of the dependence on the time t of the quantities s, dsldt and d2sldt2for harmonic varia­ tions and with <p0= 0 are shown in Fig. 28.1.

\subsection{}\label{28.1.5} It is evident from the second equation in Sect. 28.1.4 that a harmonically varying quantity s satisfies the differential equation
The general solution of this equation is of the form s= sina>t+ A2costot,
where Ai and A2 are arbitrary constants of integration. The values of Axand A2can be found from the initial conditions, i.e. when the values of s and dsldt are known for the initial in­ stant of time It = 0). Thus


The general solution can be transformed to the standard form for harmonic oscillations (vibrations) given in Sect. 28.1.3:
s= Asin(CD*+ q)0),
where A= y and cp0= arctan(A2/A1).
Thus, the quantity s has harmonic variations if and only if it satisfies the differential equation given at the beginning of this subsection. This equation
is called the differential equation
of harmonic variations (or vibra­
tions or oscillations).

\subsection{}\label{28.1.6} Harmonic oscillations (vi­
brations) can be represented
graphically by means of a vec­
tor in a plane. For this purpose
vector A is drawn from the ori­
gin of coordinates 0 in the plane
XY (Fig. 28.2). The magnitude
of vector A is equal to the am­
plitude A of the oscillation
being considered and the vector
makes the angle cp = cd* -f- cd0,
equal to the phase of the oscilla­
tion at the given instant of time t, with the coordinate axis OX. In the course of time angle <p increases so that vector A rotates at a constant angular velocity about point 0 equal to the angular frequency of oscillation (vibration) co. The corresponding pro­ jection of vector A on the vertical axis OY varies harmonically according to the relation
Ay= s= Asin(CD*+ <p0).
The graphical representation of harmonic vibrations (or oscil­ lations or variations) by means of a rotating amplitude vector is called the vector diagram method. It is extensively applied, for instance, in adding together harmonic oscillations having the same direction (28.4.2).

\subsection{}\label{28.1.7} According to the Euler formula for complex numbers,
eicp= cosq)+ isincp,
where i = V-A. Hence the harmonic oscillation (vibration) s= Asin(co*+ cp0)= Acos(m+ <px),where<px= (p0—n/2,
can be written in the exponential form
7 Aei(*t= Aei ')t+ \
where A = Aeiwi is the complex amplitude. Only the real part
of the complex function s, denoted by Re  s , has a physical meaning:
Re s = 5= Acos(co*+ <pj)= Asin(co*cp0), where <Po==91+ art/2.

\section{Mechanical Harmonic Vibrations}
\label{sec-28.2} 

\subsection{}\label{28.2.1} If a particle vibrates harmonically in a straight line along the coordinate axis OX about the equilibrium position, taken as the origin of coordinates, the dependence of the coor­ dinate x of the particle on the time * is of the form given in Sect. 28.1.3 (with s = x). Thus
x = A sin (co* -|- cp0).
The projections of the velocity v and acceleration a of the par­ ticle on the OX axis are:
vx— cos(W+ To) and ax= —aosin(co*+ cp0),
where v0= Aco is the amplitude of the velocity, and a0= = A 2— Vo<£is the amplitude of the acceleration.
The force F acting on the particle equals
F= mu and Fx= —muPx,
where m is the mass of the particle. Consequently, force F is proportional to the displacement of the particle from the equi­ librium position and opposite in direction:
F=—
where i is a unit vector of the OX axis.
Such a dependence of the force on the displacement is typical for an elastic force (3.3.6). Hence, forces of other physical ori­ gin, satisfying the same kind of relationship, are said to be quasi-elastic forces.

\subsection{}\label{28.2.2} The kinetic energy of a particle vibrating harmonically in a straight line is equal to
E*=— =-
-COS2(oJ+ CPo mcoM2 cos2( £-(-cp0),
m c,-)2A2
£k= ---t---[1+ cos (2 f+ 2<p0)].
The kinetic energy of the particle varies periodically from 0 to m  2A2/2 as it vibrates harmonically with the angular fre­ quency 2co and the amplitude mco2i42/4 about its average value of mu)2A2/4.
The potential energy of a particle vibrating harmonically in a 8raight line due to the effect of a quasi-elastic force is equal to
f r , mco2 2 mcoM2 . . = — Fxdx=----— =-- — sin2(o>*+cpo)i
0 or
A2
Ep=     -[1—cos(2 *+2<p0)]
:= 4 [1+cos(2 f+2cp0+Ji)].
The potential energy of the particle varies periodically from 0 to m  2A2/2 as it vibrates harmonically with the angular fre­ quency 2o) and the amplitude mcoM2/4 about its average value of muPAVA. The variations of the potential and kinetic energies have a phase difference of jt, so that the total mechanical energy of the particle remains constant in vibration:

E=Ek~(-Ep=---------£ =const.
Curves of the dependence of Ek, Ep and E on the time t for the case when cp0 = 0 are shown in Fig. 28.3.

\subsection{}\label{28.2.3} Example 1. A linear harmonic oscillator is a particle of mass m having harmonic vibration along a straight line by the action of elastic force Fei = — kxi (3.3.6). Serving as an example of such a system is the spring pendulum, where a weight of mass m is suspended by an absolutely elastic spring (k is a
factor specifying the elastic properties of the spring). The mo­ tion equation is
md2x kxord*xJc x=0. dt2 dt2 m
It follows from Sect. 28.1.5 that an oscillator (spring pendulum) has harmonic vibration according to the equation x =
= A sin (at + cp0) with an angular frequency co and period T equal to
“ - / J L ,»d r = 2 » / i .
The potential energy of a linear harmonic oscillator is

\subsection{}\label{28.2.4} Example 2. A compound, or physical, pendulum is a rigid body capable of oscillating under the action of the grav­ ity force mg about a horizontal axis 0, which does not pass through the centre of gravity of the body, and is called the axis of oscillation, or swing (Fig. 28.4). The centre of gravity of the pendulum coincides with its centre of mass C (7.3.2). Point 0, the intersection of the axis of
oscillation with a vertical plane passing through the centre of gravity of the pendulum and per­ pendicular to its axis of oscilla­ tion, is called its point of suspen­ sion.
In the absence of friction forces in the suspension, the motion equa­ tion of the pendulum is of the Form (4.3.4)
r d2<X v • Iat*1 —mgdsincc,
where a is the angle of devia­ tion of the pendulum about its axis of oscillation from the equi­
librium position, d — OC is the
distance from the .centre of mass
of the pendulum to its axis of
oscillation, I is the moment of inertia (4.2.1) about the same axis, m is the mass of the pendulum, and g is the free fall acceleration. Upon small oscillations of the pendulum sin a   a and the motion equation of the pendulum is of the form
d2a , mgd ( n , r~c;== ’
1.e. angle a satisfies the differential equation of harmonic vi­ brations (28.1.5). Ilence, in the absence of friction, small oscillations of a compound pendulum are harmonic:
a—a0sin(to*+ cp0),
where a0is the amplitude of variation of angle a, and
0>=  / a n d T 2nl/"—L. fI Vmgd
are the angular frequency and period of small oscillations of a compound pendulum.

\subsection{}\label{28.2.5} Example 3. A simple pendulum is a particle suspended by a weightless, inextensible thread and oscillating in a vertical plane by the action of the gravity force. The simple pendulum is the limiting case of the compound (physical) pendulum, whose whole mass is concentrated at the centre of mass, so that d — I, the length of the simple pendulum. The moment of inertia of a simple pendulum with respect to its axis of oscillation is / — ml2. Correspondingly, the angular frequency and period of small oscillations of this pendulum are
co= —andT=2n
Small oscillations of compound and simple pendulums are examples of isochronous vibrations, i.e. vibrations whose fre­ quency and period are independent of their amplitude.
In the general case, the period of a compound pendulum depends upon its amplitude a0. Thus
The variation in the value of T when aQis increased to 15°does not exceed 0.5 per cent.

\subsection{}\label{28.2.6} The equivalent length Zeq of a compound pendulum is the length of a simple pendulum that has the same period of oscillation. Thus
where I c is the moment of inertia of the compound pendulum with respect to an axis passing through the centre of mass C of the pendulum and parallel to its axis of oscillation. Point O], lying on the straight line OC at the distance of Zeq from the point of suspension O (Fig. 28.4), is called the centre of oscilla­ tion of the compound pendulum,. The centre of oscillation 01 and the point of suspension 0 are interchangeable: if the pen­ dulum is suspended so that its axis of oscillation passes through point Oi, then point 0 will coincide with the new position of the centre of oscillation, i.e. the equivalent length and period of oscillation of the pendulum remain unchanged.

\subsection{}\label{28.2.7} Example 4. Small free vibrations of electrons in a plasma (22.6.1) are called Langmuir oscillations of the plasma. They are caused by the forces of an elec­
tric field, which is set up in an elec­
trically neutral plasma as the result
of some random deviation of the
spatial distribution of the electrons
from their equilibrium distribution.
Assume that in a flat layer of plasma
of thickness Z (Fig. 28.5) the elec­
trons are displaced over the small
distance s along the positive direc­
tion of axis OX. This leads to the
formation of a surplus positive charge
at the left-hand part of the layer and
a surplus negative charge at the
right-hand part. Consequently, an electric field is set up, whose strength E is along the OX axis and its projection on this axis is E x = n0es/s0l where n0 is the electron concentration in the plasma, e is the magnitude of the charge of the electron, and e0 is the electric constant. According to Newton’s second law, the motion equation for the electrons of the plasma in this elec­ tric field is of the form:
d*s—eEx— e2n0sord*s e2n0s=Q, dta dt2 me0
where m is the mass of the electron. Thus, the electrons of the plasma have free harmonic vibrations at the angular frequency


which is called the Langmuir plasma frequency, or simply, the plasma frequency.

\section{Free Harmonic Oscillations in an Oscillatory Electric Circuit}
\label{sec-28.3} 


\subsection{}\label{28.3.1} An example of an electric circuit in which free electric oscillations may occur is the simple oscillatory circuit (Fig. 28.6). It consists of a capacitor of capacitance C connected in series to
I>0
Fig. 28.6
an induction coil with the self-induc- tancc L. When the previously charged capacitor is connected to the coil, free oscillations of the capacitor charge and the current in the coil are pro­ duced. This sets up a variable elec­ tromagnetic field that propagates through space at a velocity equal to that of light. Hence, if the linear di­ mensions I of the circuit are not too large (Z< c/v, where c = 3 X
Xl08m/sis the velocity of light in vacuum, and v is the frequency of oscillation in the circuit), it can be assumed that at each instant of time t the current I is the same at all parts of the circuit. Such a variable current is
said to be quasi-stationary. According to Ohm’s law (21.2.3) for the portion l-L-2 of the circuit (Fig. 28.6)
/i?=F1-F2+les or i-J.


Here q and Vx — V2 = — qtC are the charge of the capacitor and the potential difference of its plates at the arbitrary instant t of time being considered, R is the electric resistance of the oscillatory circuit, i.e. of the part l-L-2 of the circuit, and
= —L (dl/dt) is emf of self-induction in the coil (25.2.4). If follows from the law of conservation of electric charge (14.1.3) that the quasi-stationary current in the circuit is I — dqldt. Therefore, the differential equation for the variation of charge q is of the form
d2q R dq , q dt2 L dt LC
0.

\subsection{}\label{28.3.2} Free electric oscillations in an oscillatory circuit are harmonic if its electrical resistance R — 0. Thus
The angular frequency go and the period T of these oscillations satisfy the Thomson formula:
( 0            = a n d T = 2 n Y L C . V LC
The charge q of the capacitor and current I in the circuit vary according to the equations
<1 -=qQsin (otf+ epo) and
/ -/„ cos (<oi+ ) = /„ Sin (co*+ (p0+ - -j ,
where q0 is the amplitude of the capacitor charge, I 0 — cog0 =
-qj YLC isthe amplitudeof the current, and cp0is the ini­ tial phase of capacitor charge oscillation. The current in the circuit has a phase lead of ji/2 over the capacitor charge.
The potential difference u = V2— V1 of the capacitor plates also varies according to a harmonic law and is in phase with charge q. Thus
u= - -=U0sin (cot-j-tp0),
where U0— qJC is the amplitude of the potential difference. The amplitude of the current is
'o = Uo ~jrf •
The quantity YL/C is called the characteristic, or wave, impe­ dance of the oscillatory circuit.

\subsection{}\label{28.3.3} With free harmonic oscillations in an oscillatory circuit, there is a periodic transfer of the energy E e of the electric field

of the capacitor into energy Em of the magnetic field of the in­ ductance coil and back again. Thus
Ee=~£c~=lxTSin*(<0,+ (f’o)= '!c~ f1—cos(2w*+2<p0)],
Em =
T T2 r T2 T T2
cos2((at+ cp0)= — [1+ cos (2 f+ 2<p0)].
For this reason, the oscillations in an electrical oscillatory cir­ cuit are frequently said to be electromagnetic oscillations in the circuit.
The values of Ee and Em vary in electromagnetic oscillations within the limits from 0 to their maximum values, equal, re­ spectively, to q/2C and L/J/2, in which q/2C = LI/2. The variations of E e and Em are displaced in phase. At the instants of time when Ee= 0, Em= Emmax= LIq/2. On the other hand,whenEm—0,Ee= Eemax= q2C.Thetotalenergy of electromagnetic oscillations is constant with time. Thus
E= Ee+ Em= =const.

\section{Adding Harmonic Oscillations}
\label{sec-28.4} 

\subsection{}\label{28.4.1} The addition of oscillations (vibrations or variations) consists in finding the equation for the resulting oscillation of a system when this system simultaneously participates in several oscillatory processes. Distinction is made between two limiting cases: the combining of oscillations having the same direction and the combining of mutually perpendicular oscillations. The first case corresponds, for example, to the vibration of weight 1 (Fig. 28.7) which vibrates with respect to weight 2 on spring a, and together with weight 2 on spring b. This same situation occurs in superposing the variations of scalar physical character­ istics of an oscillatory or variable system (pressure, temperature, density, electric charge, current, etc.).

\subsection{}\label{28.4.2} Two harmonic oscillations having the same direction, such as =A1sin(co -|-(p2)ands2=A2sin(co2£+(p2)canbe added together by using the vector diagram method (28.1.6),


Shown in Fig. 28.8 are the vectors Ax (t) and A2 (t) of the first and second oscillations at the arbitrary instant t of time when the phases of these oscillations equal: (/) = coxf -|- (px and d>2(0 = + 92- The resultant oscillation s ~ sx+ s2 cor-
iesponds to the vector A(i)= Ax(t)+ A2(l) whose projec­ tion on the vertical axis OY is equal to s. Thus
.v= A (t) sin O (t). According to the cosine law
[A(t)]2= A*+ A*+ 2A,A2c o s[<Da(I)- ipj(01, and
tanO(t)= AjsinOx(t)-j-AzsinP2(0
AxcosOx(t)+A2c o s02(t)

\subsection{}\label{28.4.3} Two harmonic oscillations sx and s2 are said to be coherent if their phase difference is independent of time. Thus
-j-[d>2(0— (01==0 and 2(0—d>i(£)=const.
Since (t>2 (t) — (t) — ((o2 — (iq) t + (92 — <Pi), the angular frequency of coherent oscillations should be the same, i.e. (o2= co1= o). At any instant of time the phase difference of coherent oscillations is equal to the difference of their initial phases:   2(t)— Oj (t)= cp2— cpj. Consequently, the resultant oscillation is harmonic and has the same angular frequency co, i.e.
s= +s2=Asin(co£+(p0),
where
A2= A2+ A 2AXA2cos(tp2—cp,) and
tan(Po=  »i"q>,+ sinq,a
Axcos cpj-j-A2cos cp2

In accordance with the values of the initial phase differences of ihe oscillations being added together, the amplitude A of the resultant oscillation varies within the limits
fromA= |i4x—42|at<p2—<Pi= it(2m+ 1)n toA= Ai-A2 atcp2—91= ±2mJt
where m = 0, 1, 2, ... is any integral nonnegative number. When cp2— 9i = d=2mjt, the oscillations being added are said to be cophased (or simply, in phase), whereas at cp2—cpx= = ±(2m+ 1)jitheyaresaidtobeaniiphased(orinopposition). 


\subsection{}\label{28.4.4} Harmonic oscillations of different frequency (co2=£ (Oj) are incoherent because their phase difference, equal to (co2 — (Oj) t + (tp2 — (pi), varies continuously with time. If such oscillations are superposed, an anharmonic resultant oscillation is obtained. The amplitude vectors At and A2of the added oscillations (Fig. 28.8) rotate at different velocities, so that the parallelogram constructed on the vectors is continuously deformed and its diagonal, vector A of the resultant oscillation,
varies in length and rotates at variable angular velocity. Two harmonic oscillations having different angular frequencies ocq and (o2can be approximately considered to be coherent only during the time interval At in which the phase difference of these oscillations varies only slightly: | w2 — \ At <C 2n or At <c

•C TCOh» where xCOh = 2it/| co2 —  j | is the coherence time of the oscillations being considered.

\subsection{}\label{28.4.5} Anharmonic oscillation obtained as a result of super­ posing two harmonic oscillations having the same direction and slightly differing frequencies (| <o2— I < wi) produces what are called beats. It proves convenient to take as the time reference /the instant when the phases of the two oscillations being super­ posed, s1 and Sg, coincide and are equal to cp0. Then sx  — — 4!sin(co + cp0) and s2=  42sin (u>2t+ q>0)= —A2sin[ccqz+ <p0+ 9(01»where (p(t)= (ca2—tOj)t. The resultant oscillation satisfies the relation
s — A (t) sin [co  + 9 0 + 1) (£)], where
A(t)]2= A2+ A 2A±A2cos<p(/) and
A2sin cp(t) Ax+A2c o scp(t)'
tan i|)(t)
In particular, when A1= A2= A0,
A t)=2A0cos 0)2 01 t and i|)(t)=
so that
X->A0cos— ——tsin ——H-<Poj.
I’he quantity |A (t) |, characterizing the swing of the oscilla- lions in beating, varies within the limits from At — A2I to .1, + A2 with the angular frequency Q = |co2— coj |, called Ilie angular beat frequency. Since the beat frequency is much lower than the frequency of oscillation (Q <c coj), the variable quantity |A (/) | is conditionally called the beat amplitude. Tlie beat period 2+ and the beat frequency are equal to
2jt 2n 1
Q 111
T, T,|
and
1
V|,= — ‘I)

|v., — Vi I,
where Tlt v,, 7’2and v2are the periods and frequencies of the superposed oscillations. Curves of the dependence of s on t in beating are shown in Fig. 28.9 for the case when Ai A2

\subsection{}\label{28.4.6} As a result of the superposition of harmonic oscillations, coinciding in direction and having angular frequencies that are integral multiples of their lowest frequency, i.e. co, 2(0, 3co, etc., a periodic nonharmonic oscillation with the period T = 2n/co is obtained. On the other hand, any complex periodic oscillation s   / (t) can be represented as the sum of simple harmonic oscillations having angular frequencies, which are integral multiples of the basic angular frequency co = 2nlT, where T is
the period of the oscillation. Thus oo
s=/(/)==-y*+ 2 cosn(0t+bnsin
n—1 oo
/4„ sin (»«>(+ «Pn).
+
n—1
where
T/2
(ln-~~rp~ | f(t)cosn(s)tdt (rc—0,1,2,...),
-T/2 T/2
bn=~jT   /(0sinrcco£d£“(w=l,2,...). -T/2
Such a representation of the periodic function / (£) is called the expansion of this function into a Fourier series, or the harmonic analysis of a complex periodic oscillation. The terms of the Fourier series corresponding to the harmonic oscillations with the angu­ lar frequencies co, 2co, 3co, etc. are called the first, or fundamental, second, third, etc. harmonics of the complex periodic oscillation s = / (t). The whole set of these harmonics forms the spectrum of the oscillations s = f (t). The composition of the spectrum depends on the form of the periodic function s = f (t). In the simplest cases, the spectrum can consist of a small number of harmonics.

The oscillation spectrum is often understood to be the frequency spectrum, i.e. the set of frequencies of simple harmonic oscil­ lations that can be superimposed to obtain the complex oscil­ lation being considered. Periodic oscillations have discrete (line) frequency spectra.

\subsection{}\label{28.4.7} Nonperiddic oscillations have, as a rule, continuous fre­ quency spectra, i.e. they can be represented as the result of the superposition of a great many harmonic oscillations, whose frequencies have all possible values in a certain range (from 0 to oo in the general case). The harmonic analysis of such oscil­ lations consists in presenting them in the form of a Fourier integral:
s — f ( t ) — j * [ a ( c o ) c o s c o t - \ - b ( c o ) s i n c o £] d c o ,
t((0)=A (/(g)sinco£dl. — OO

Certain nonperiodic oscillations, called almost periodic (quasi- periodic) oscillations have a line frequency spectrum. But the included angular frequencies are incommensurable with one another, because their ratios are expressed by irrational numbers. 28.4.8 Modulation is the change of any parameter of a periodic oscillation (for instance, the amplitude or frequency) according to a definite law during a length of time considerably longer than the period of the oscillation.

In amplitude modulation, for example, of the harmonic oscillation s = A0sin (cd0£-f- <p0), the modulated oscillation is of the form
s= A0[i+ b(*)]sin(cd0*+ cp0),
where|b(t)< 1.

If amplitude modulation is carried out according to the harmonic
law b(t)= b )cosQ/, where b0= const and Q<C then
s= A0(1—b0cosQt)sin(co0£+ <Po)-
This modulated oscillation has a line frequency spectrum because it can be represented as the sum of three harmonic oscillations with the angular frequencies co0, co0— Q and co0-f- Q, and amplitudes equal, respectively, to A0,A0b0l2 and A0b0/2. Thus
(1+ fy)cosQ0sin(a t+ cp0)
= A0sin (<n0<+ (f-0)+ - | - sin |(co0+ £2)<+ q)0]
+ sin [(co0—Q)<+ cp0] .

In frequency modulation of the harmonic oscillation s = A0 X X sin (cd0£-f- cp0), carried out according to a harmonic law, the
modulated oscillation is of the form s-=A0sin[cd0(1+ b0cosQ/)t+ (p0],
whereb0< 1andQ< co0.

Accordingly, in phase modulation, the initial phase of the oscillation is changed. Thus
s= A0sin((o0/,+ AcpcosQ/), where Q< co0.
In the general case oscillations can be modulated simultaneously in amplitude and in phase (or frequency). An example of such modulated oscillations is beating (28.4.4).

\subsection{}\label{28.4.9} Addition of mutually perpendicular harmonic oscillations nj equal frequency. Assume that point M oscillates simultaneous­ ly along the coordinate axes OX and OY according to the law: x = Axsin (rn + tpj) and y = A2sin (o)t + q)2), where x and / are the Cartesian coordinates of point M. The path equation lor the resultant motion of point M in plane XOY can be found by eliminating parameter t from the expression for x and y. Thus
2xy COS(92—cpi)= sin2(t2—9i)- A\ AxA2
This path is an ellipse (Fig. 28.10) and point M describes this ellipse during a time equal to the period of the oscillations being

added together (T = 2jx/co). Hence, the resultant motion of point M is called elliptically polarized oscillation.

The orientation of the axes of the ellipse in plane XOY, as well as its dimensions depend upon the amplitudes A1and A2of the oscillations being combined and on their initial phase difference 92 — 9i- If 92 — 9i = (2/n + 1) Jt/2, where m = 0, ±1, ±2, . . the axes of the ellipse coincide with the coordinate axes OX and OY, and the lengths of the semiaxes are equal to the amplitudes Ax and A2. Thus

II, in addition, A1— A2, the path of point M is a circle. Such resultant motion of point M is called circularly polarized oscil­ lation.

When<p2—cpi=mn,wherem=0,+1,±2,.. theellipse degenerates into a straight line. Thus
The plus sign corresponds to even values of m, i.e. to the addition of cophased oscillations (Fig. 28.11a); the minus sign, to odd values of m, i.e. to the addition of oscillations that arc in
opposition (antiphased) (Fig. 28.11 ). In either case, point M executes linearly polarized oscillation. It oscillates with harmonic motion at the frequency of the oscillations being combined and
with the amplitude A = YA\ + A\ along a straight line that makes the angle a = arc tan ( 42c°s rn.nlAx) with the axis OX. 28.4.10 In the addition of mutually perpendicular oscillations having the angular frequencies pw and goo, where p and q are whole numbers, we obtain
x — A1sin (p(0t + cpj) and y = A2sin (q(ot + cp2).
The values of coordinates x and y of the oscillating point M are simultaneously repeated after equal periods Tn of time, equal
to the least common multiple of T1— 2nlp(o and T2— 2ji/</(jd, which arc the periods of oscillation along the OX and OY axes. Hence, the path of point M is a closed curve whose shape depends upon the ratios of the amplitudes, frequencies and initial phases of the oscillations being combined. Such closed paths of point M, which has simultaneous harmonic oscillations in two mutually perpendicular directions, are called Lissajous figures.


Lissajous figures arc inscribed into rectangles whose centre coin­ cides with the origin of coordinates and whose sides are parallel to coordinate axes OX and OY. The distances of these sides from the axes are equal, respectively, to A2and Ax. The ratio of the frequencies qco and pco of the combined oscillations is equal to the ratio of the number of times the corresponding Lissajous figure touches a side of the rectangle parallel to the OX axis to the number of times it touches a side parallel to the OY axis. Shown in Fig. 28.12 are the Lissajous figures for three different values of the ratio qlp (2:1, 3:2 and 4:3) and for initial phase differences of Acp—<px—cp2= jt/2.



